\documentclass{article}
\usepackage{graphicx} % Required for inserting images
\usepackage{datetime} % Add this line to include the datetime package
\usepackage{amsmath}
\usepackage{amssymb}
\usepackage{amsthm}
\usepackage{tikz}
%\usepackage{hyperref}
\usepackage{xcolor} % Required for defining custom colors

\definecolor{darkblue}{rgb}{0.0, 0.0, 0.5} % Define a darker blue color

\usepackage[colorlinks=true, linkcolor=darkblue, urlcolor=darkblue, citecolor=darkblue]{hyperref} % Use the darker blue color

% Define theorem style
\theoremstyle{plain} % Bold title, italic body
\newtheorem{theorem}{Theorem}
\newtheorem{lemma}[theorem]{Lemma}
\newtheorem{corollary}[theorem]{Corollary}

\theoremstyle{definition} % Bold title, normal body
\newtheorem{definition}[theorem]{Definition}
\newtheorem{example}[theorem]{Example}


\title{Homework 1}
\author{Aaron Honculada}
\date{\monthname[\month] \the\year}

\begin{document}

\maketitle

\section*{Problem 1}
\noindent\textit{Determine which of the following are groups.}\\
\noindent\text{(a)} The set of integers $\mathbb{Z}$ under the operation $x \star y = x + y + 1$.

We check the group axioms:
\begin{itemize}
    \item \textbf{Closure:} Let $x,y \in \mathbb{Z}$. Then $x \star y = x + y + 1 \in \mathbb{Z}$.
    \item \textbf{Associativity:} Let $x,y,z \in \mathbb{Z}$. Then $(x \star y) \star z = (x + y + 1) \star z = (x + y + 1) + z + 1 = 
        x + (y + z + 1) + 1 = x \star (y + z + 1) = x \star (y \star z)$.
    \item \textbf{Identity:} Let $x \in \mathbb{Z}$ and $e = -1$. Then $x \star e = x + -1 + 1 = x + 0 = x$. Similarly,
    $e \star x = -1 + x + 1 = x$.
    \item \textbf{Inverse:} Let $x \in \mathbb{Z}$ and $y = -x-2$. Then $x \star y = x + y + 1 = x + (-x-2) + 1 = -1 = e$.
        Similarly, $y \star x = -1 = e$.
\end{itemize}
Therefore, $\mathbb{Z}$ under the operation $x \star y = x + y + 1$ is a group.
\\

\noindent\text{(b)} The set of natural numbers $\mathbb{N}$ under the operation $x \star y = \text{max}(x,y)$.

First we identify the identity element. Let $e = 1$. Then $x \star e = \text{max}(x,1) = x$ for all $x \in \mathbb{N}$.
Similarly, $e \star x = x$.
However we see that for any element in the natural numbers, there is no inverse element $y \in \mathbb{N}$ such that $x \star y = e$.
Therefore, $\mathbb{N}$ under the operation $x \star y = \text{max}(x,y)$ is not a group.
\\

\noindent\text{(c)} The set of rational numbers $\mathbb{Q}$ under the operation $x \star y = xy + x + y$.

We show this operation is not associative. 
\begin{align*}
    (x \star y) \star z = (xy + x + y) \star z = (xy + x + y)z + (xy + x + y) + z = xyz + xz + yz + x + y + z
\end{align*}
\begin{align*}
    x \star (y \star z) = x \star (yz + y + z) = x(yz + y + z) + x + (yz + y + z) = xyz + xy + xz + x + y + z
\end{align*}

We see that $(x \star y) \star z \neq x \star (y \star z)$ for all $x,y,z \in \mathbb{Q}$.
Therefore, $\mathbb{Q}$ under the operation $x \star y = xy + x + y$ is not a group.
\\


\section*{Problem 2}
\noindent\textit{Fix a positive integer $n\geq 3$, and let $\omega = e^{2\pi i / n}$.
Define functions $r_a, s_a : \mathbb{C} \to \mathbb{C}$ for $a = 0,\ldots,n-1$
by 
\begin{align*}
    r_a(z) = \omega^a z, \quad s_a(z) = \omega^a\overline{z},
\end{align*}
where $\overline{z}$ is the complex conjugate of $z$.}

\noindent\text{(a)} Show that $\{r_0,\ldots,r_{n-1},s_0,\ldots,s_{n-1}\}$ forms a group under composition.

We show that this set forms a group under composition.
\begin{itemize}
    \item \textbf{Closure:} Let $r_a, r_b \in \{r_0,\ldots,r_{n-1},s_0,\ldots,s_{n-1}\}$. Then $r_a \circ r_b(z) = r_{a+b}(z)$.
        If $a+b \geq n$, then $r_{a+b}(z) = r_{a+b-n}(z)$.
        Similarly, $s_a \circ s_b(z) = r_{a-b}(z)$. This is because
        \begin{align*}
            s_a \circ s_b = \omega^a \overline{\omega^b \overline{z}} = \omega^a \overline{\omega^b} \overline{z} = \omega^{a-b} \overline{z} = r_{a-b}(z)
        \end{align*}
        Composing $s_a$ with $r_b$ results in $s_a \circ r_b(z) = \omega^a \overline{\omega^b z} = \omega^{a-b} \overline{z} = r_{a-b}(z)$.
        We have shown closure under composition.
    \item \textbf{Associativity:} Composition of functions is associative. For instance,
        \begin{align*}
            (r_a \circ r_b) \circ r_c = r_{a+b} \circ r_c = r_{a+b+c}
        \end{align*}
        \begin{align*}
            r_a \circ (r_b \circ r_c) = r_a \circ r_{b+c} = r_{a+b+c}
        \end{align*}
        Similarly,
        \begin{align*}
            (s_a \circ s_b) \circ s_c = r_{a-b} \circ s_c = r_{a-b-c}
        \end{align*}
        \begin{align*}
            s_a \circ (s_b \circ s_c) = s_a \circ r_{b-c} = r_{a-b-c}
        \end{align*}
        We have shown associativity.
    \item \textbf{Identity:} The identity element is $r_0(z)$.
    \item \textbf{Inverse:} Given that there are a finite number of $n$ unique functions $r_a$ then the inverse of
        each $r_a$ is some $r_b$ where $a+b = n$. The inverse of $s_a$ is $s_a$ since $s_a \circ s_a = \omega^{a}\overline{\omega^{a}} = 1 = \omega^0 = r_0$.
\end{itemize}

\noindent\text{(b)} Explain why this group is isomorphic to the dihedral group
$D_n$, the group of symmetries of a regular $n$-gon.

This group is isomorphic to $D_n$ because it has the same structure since
we can map each element in $D_n$ that is a rotation to $r_a$ and each element in $D_n$ that is a reflection to $s_a$.


\section*{Problem 3}
\noindent\textit{Suppose $G$ is a group in which every element $x\in
G$ satisfies $x^2 = e$. Show that $G$ is abelian.}

Given that $G$ is a group and every element $x \in G$ satisfies $x^2 = e$, we have that
for all $a,b \in G$ we know $ab \in G$ and every element $a$ is its own inverse $a^{-1} = a$. Thus we have that
\begin{align*}
    (ab)^2 = e \implies abab = e \implies ab = (ab)^{-1} = b^{-1}a^{-1} = ba
\end{align*}
Therefore $ab = ba$ and $G$ is abelian.

\section*{Problem 4}
\noindent\textit{Let $G$ be a finite group with an even number of elements.
Show that $G$ has an element of order 2.}

In $G$ we know that the identity element $e$ has order 1. Given the order of $G$ is even, we can call the
order of $G$ to be $2n$ for some $n \in \mathbb{N}$. Without the identity element, we have $2n-1$ elements left.
We know that every element in $G$ has a unique inverse. Thus we can pair each element $a \in G$ with $a^{-1}$.
At most we can pair $2n-2$ elements. We call the only unpaired element $b$. Since $b$ is unpaired,
we conclude it is its own inverse $b = b^{-1}$. Thus $b^2 = e$ and $b$ has order 2. If less than $2n-2$ elements 
can be paired, then the remaining elements must be paired with themselves. Thus they are their own inverses and have order 2.

\section*{Problem 5}
\noindent\textit{Let $G$ be a nonempty set closed under an associative operation $\star$.
Suppose that for all $x,y,x \in G$, if $x \star z = y \star z \text{ or } z \star x = z \star y$, then $x = y$.}

\noindent\text{(a)} Prove that if $G$ is finite, then it is a group.
\begin{proof}
    Given that $G$ with the associative operation $\star$ is closed under the operation,
    we must show that $G$ has an identity element and that every element has an inverse.

    It follows from $G$ having the left and right cancellation properties that $\star$ is injective.
    Since $G$ has finite elements, is closed under $\star$, and no two elements are mapped to the same element, $\star$ is also surjective.
    
    Therefore, $\star$ is both injective and surjective and thus a bijection.
    Since we can map each element in $G$ to an element in $G$, $\star$ is a bijection,
    and thus each element in $G$ has an inverse.

    Since $G$ is closed under $\star$, we know that the identity element $e$ is in $G$.
    Therefore, $G$ has an identity element and each element has an inverse.
    Thus $G$ is a group.

\end{proof}

\noindent\text{(b)} If $G$ is infinite, must it be a group?

If $G$ is a set of infinite elements, then we can no longer claim that
$\star$ is surjective. This means that the cancellation property does not 
necessarily imply that $\star$ is a bijection, and we cannot conclude that $G$ is a group.

\end{document}